\section{Le prototype Codex: architecture et choix techniques}
Le prototype interne \enquote{Codex}, développé sous le framework open source Python-Django par le responsable du Département d'Information Technique (DIT) de la bibliothèque Cujas, s'inscrit dans une rupture avec le système historique XTF. Ce dernier, bien que bénéficiant d'une solide réputation dans le monde anglo-saxon, est devenu aujourd'hui obsolète due à l'absence de mise à jour depuis 2018. Cette situation engendre d'importants risque de sécurité, une menace de rupture de service par le réseau Renater et une déconnexion par rapport aux standards utilisées. Ce sont ces multiples raisons qui ont poussé le responsable à développer ce prototype. Codex a été réfléchi et conçu afin de sortir la bibliothèque de son isolement fonctionnel.

\subsection{Le dépassement de l'arborescence \enquote{à plat} et la modélisation hiérarchique}
L'un des obstacles majeurs identifiés lors de l'audit de la première bibliothèque numérique, CujasNum réside dans la rigidité architecturelle développé sous XTF. Cette infrastructure limitait l'implémentation d'autres types de documents nécessitant une arborescence complexe. En effet, comme le souligne Isabelle Westeel, \enquote{toute gestion de collections numériques impose de tenir compte de la granularité (physique ou logique)}. La structure \enquote{à plat} du logiciel a donc imposer à la bibliothèque Cujas la diffusion de monographies uniquement sur CujasNum. 

Pour le Département des périodiques et des ressources électroniques (DPRE), cette contrainte technique empêchait l'intégration de revues juridiques. Un périodique est par nature ce que l'on qualifie \enquote{d'un objet composé (compound object)}, c'est-à-dire un \enquote{assemblage de texte}, dont l'arborescence doit être préservée pour ne pas en perdre l'intelligibilité scientifique. Une revue est composé différent d'une monographie. Elle se compose alors \enquote{d'un tire, de l'année de publication, du volume} au minimum, mais aujourd'hui le DPRE requiert une granularité jusqu'à l'article lui-même. 

La souplesse du prototype Codex permet selon le responsable DIT \enquote{d'introduire une modélisation de données à plusieurs niveaux}. Elle permet ainsi d'ajouter d'autres formats de documents sans avoir à \enquote{créer de faux PDF pour pouvoir les intégrer}, dans le cas où l'on avait souhaité implémenter, par exemple, le fonds iconographiques. Sous XTF, la notion de \enquote{collections} était absente : toute création de collection était impossible. Celle crée sur CujasNum relève \enquote{d'un bricolage maison: réalisé à l'aide d'une requête SQL}, permettant de les constituer. Désormais, le prototype résout cette obstacle rencontré par les experts métiers en mettant directement à leur disposition la possibilité dans le back-office de pouvoir créer des collections à l'aide de sets dynamique. Le système offre la possibilité de les constituer à l'aide de plusieurs filtres que l'expert métier peut lui-même consituter : par sujet ou indexation thématiqque, par cote ou classification Dewey, par couverture spatiale ou géographique ou par lot de numérisation ou par critère temporel. Sur le plan technique, cette agilité repose sur la transition du Dublin Core simple vers le Dublin Core qualifié, c'est-à-dire \enquote{des élements de raffinements peuvent compléter certains éléments}, remplaçant \enquote{les éléments de base par des éléments plus précis}. \enquote{Ainsi, au lieu de title (dc:title), on pourra utiliser alternative (dcterms: alternative) pour désigner un titre abrégé ou un titre traduit}. Cela permet aux usagers, une navigation plus fluide. Ce prototype offre la possibilité de disposer d'une navigation ergonomique et d'une indexation fine plus agréable pour les usagers de la bibliothèque numérique, ainsi qu'une promesse de confort pour les experts métiers dans le back-office.

Au delà, du confort que cela apporte sur la page d'accueil de la bibliothèque numérique, l'implémentation de ces sets revêt un intérêt stratégique majeur pour l'interopérabilité de la BIU Cujas. En effet, l'activation d'un entrepôt OAI-PMH permet à Codex d'exposer ces sets de manière structurée, favorisant un moissonnage sélectif par les autres institutions patrimoniales, comme la BNF. Comme le souligne Vincent de Lavenne, le moissonnage sélectif permet \enquote{aux fournisseurs de service de sélectionner les données qu'ils moissonnent, et donc de proposer un service pertinent}, évitant le stockage inutile de l'information et \enquote{de construire un accès adapté aux archives ouvertes}. 

\subsection{La déconnexion du traitement des images par le standard IIIF}
Sous CujasNum, pour consulter des ouvrage volumineux, les usagers n'avaient d'autres choix de télécharger des fichiers PDF extrêmement lourd, pouvant atteindre \enquote{des tailles de 300 à 400 go}. Une telle concentration de données provoquait des saturations systématiques de la bande passante et \enquote{des échecs de téléchargements récurrents}, générant une frustration chez les chercheurs et un abandon des étudiants face à ces PDF massifs. 

Pour surmonter cette limite de performance, le prototype Codex opère une transition architecturale en séparant la gestion des données de celle des images. Il s'appuie sur le serveur IIIF indépendant. Ce protocole, devenu le standard incontournable des grandes institutions patrimoniales françaises et internationales, fonctionne de manière découplée : entre la distribution des images à la demande et la génération de manifestes JSON. Ce manifeste est produit automatique sous le format JSON-LD. Ce fichier de métadonnées, extrêmement léger, décrit la structure intellectuelle du document avec l'ordre des pages et l'emplacement physique des images, permettant ainsi aux chercheurs d'exploiter les données brutes et pouvant procéder à des extractions d'informations à l'aide de script.

Grâce au serveur IIIF, le prototype Codex propose l'intégration de deux visionneuses : la visionneuse Turner et celle plus avancée Mirador. Ce choix, peu commun, offrirait à la bibliothèque Cujas la possibilité de concevoir une interface de consultation flexible et adaptée à différents profils d'utilisateurs : entre le grand public et le chercheur. Lors du benchmark réalisé pour voir comment les bibliothèques avaient réfléchi à la consultation de leurs documents, ce mode de fonctionnement a été observé uniquement dans le monde anglo-saxon, sur la Digital Bodleian Library d'Oxford et sur celle de Cambridge. Une attention doit porter par la bibliothèque Cujas face à ce choix, peu courant, car pensant s'adapter au mieux aux attentes des usagers, il ne faudrait pas que cela les rendent confus face à la présence de deux visionneuses. En effet, il semble que cela ne soit pas une habitude perpétrée par les institutions patrimoniales françaises car celle-ci propose uniquement une visionneuse IIIF. Certains ont choisi la visionneuse Universal Viewer tandis que d'autres utilisent Mirador, répondant ainsi à des objectifs différents et des attentes propres à leurs usagers.

\subsection{L'interopérabilité des données et la pérennité des identifiants (ARK)}
Pour s'intégrer au web sémantique et aux grands réservoirs nationaux, comme le projet Orion de l'Abes, le prototype s'appuie exclusivement sur des standards ouverts. La bibliothèque doit proposer des services permettant l'utilisation brute de données ainsi que des services d'export s'intégrant naturellement dans les chaînes de travail personnelles des chercheurs. Cette appropriation de la donnée par les chercheurs se heurte aujourd'hui à la volatilité des sessions de recherche. Dans CujasNum, l'absence de compte utilisateur permanent condamnait le chercheur à voir son panier de documents sélectionnés disparaître dès la fermeture du navigateur. Ce mode de fonctionnement oblige le chercheur à refaire inlassablement les mêmes requêtes d'une session à l'autre. Pour surmonter cet obstacle, la théorisation de l'espace personnel par Elina Leblanc s'avère la solution idéale. Leblanc préconise la création d'un espace personnel conçu comme un véritable \enquote{carnet de lecture}. 

Pour cela, la bibliothèque Cujas doit transformer un stock de fichiers en un environnement de connaissances éditorialisées. Aujourd'hui, le rôle des bibliothèques ne se limite pas à stocker des fichiers isolées; il s'agit d'assembler, de contextualiser et de faire dialoguer ces ressources entre elles pour en faire émerger une plus-value intellectuelle et de nouvelles connaissances. Dans leur ouvrage Antoine Isaac, Emmanuelle Bermès et Gautier Poupeau expliquent ce phénomène à travers le concept de la \enquote{citation de contenus}. Pour les auteurs, elle consistait à \enquote{sélectionner des ressources sur le Web, de les réorganiser, de les annoter et de donner à consulter ce travail dans une interface particulière permettant de saisir le sens qui émerge de la juxtaposition de ces ressources}. Dans ce troisième âge, ce traitement intellectuel est l'un des moyens utilisés par les institutions patrimoniales transformant le document-objet en un espace d'apprentissage  et de valorisation du travail scientifique réalisés par des spécialistes.

Pour l'établissement, l'enjeu de cette refonte doit permettre de développer leur prototype Codex, répondant aux attentes internes et externes. Le non-recours à une plateforme open-source de gestion de contenu (CMS), à l'instar d'Omeka S, ne doit pas contribuer à son isolement des autres institutions patrimoniales. Face à cet obstacle, la bibliothèque Cujas doit se rapprocher des standards s'intégrant au Web de données. Dans l'Enseignement Supérieur et de la Recherche (ESR), le logiciel libre Omeka S a été largement utilisé par les institutions pour la création de leur bibliothèque numérique. Malgré l'aspect \enquote{clé en main} d'Omeka S, la maintenance de ce type de système peut s'avérer lourde, demandant de posséder des compétenteces informatiques permettant d'adapter les mises à jour logiciel. L'aspect communautaire de ce logiciel permet à ceux n'ayant pas de service informatique dédié de s'entraider face à des problèmes communs. Malgré les inconvénients que ce type de développement peut représenter n'étant jamais complètement adapté aux particularités propres de l'établissement développant sur-mesure représente un risque fort pour l'institution. Pour surmonter, ce dilemme le prototype se doit d'utiliser des formats ouverts et standardisés permettant aux institutions d'utiliser leurs données, évitement ainsi l'isolement technique. 

Le choix de s'appuyer sur des formats universels et des API vise à garantir la visibilité, la réutilisation et de la diffusion des collections à l'échelle nationale. Contrairement à l'entrepôt OAI-PMH de CujasNum, devenu difficile à moissonner voir impossible, le prototype Codex permet de réactiver le moissonnage en le standardisant davantage. Le système intègre dorénavant des API REST capables de délivrer des données structurées et interopérables sous formats XML ou JSON, rompant avec l'image en silo. Cette architecture modulaire offre une double opportunité : elle permet à la BIU Cujas de moissonner d'autres portails comblant ces collections physiques ainsi que d'exposer ses propres métadonnées pour être moissonnées, augmentant drastiquement la visibilité de ses fonds.

Une des fortes attentes des chercheurs repose sur la stabilité à long terme des citations et de l'accès aux ressources numériques, pour qui la fiabilité d'un lien hypertexte est une condition importante de la scientificité de leurs travaux. A cette fin, il envisage \enquote{de désigner ARK vers la notice bibliographique plutôt que le PDF}. A la bibliothèque Cujas, la méthode de création des identifiants pérennes est forgé à partir \enquote{du code-barre de l'ouvrage}, garantissant le lien permanent évitant les doublons. Cela permettrait ainsi de lien numériquement la notice bibliographique à celle du catalogue.

\subsection{La gestion fine des droits d'accès}
La bibliothèque numérique souhaite numériser des documents sous droit à intégrer dans sa nouvelle plateforme. Tout en restant dans un cadre strictement légale, la bibliothèque se doit de réfléchir à l'accès qu'elle permet pour ces documents ne pouvant être exposé librement sans contrainte. Il est nécessaire pour l'établissement d'établir clairement une charte légale indiquant distinctement ce que l'usager a le droit de faire ou non concernant un document libre de droit et celui sous droit. La première étape consiste à assortir la notice du document à une \enquote{licence Etalab et/ou Creative Commons}. Cujas aimerait développer \enquote{un flitrage des documents sous droit ou libre de droit, grâce à la balise DcRight des notices Dublin Core}. Une des bibliothèques le fait également. Il s'agit de Polona, la bibliothèque nationale de pologne, permettant de filtrer les documents sous droit. La deuxième étape serait de mettre en place \enquote{un horodatage débloquant le document sur le web une fois que les droits sont arrivés à terme}.

Le prototype Codex permet de mettre en place des protocoles d'authentification centralisée SSO et SAML. En s'appuyant sur la fédération d'identité de Renater, grâce à ces protocoles, le système interroge à distance les serveurs des universités partenaires comme Paris I et Paris II. Cette architecture modulaire permet d'appliquer une gestion différenciée des droits d'accès. Les métadonnées descriptives restent accessible publiquement sur le Web, tandis que la consultation des images est restreinte aux droits accordés par la bibliothèque en fonction du statut occupé par la personne. Un autre moyen d'accès, plus traditionnel, est de se connecter aux postes physiques connectés au réseau intranet local au sein de Cujas.

Sous XTF, la chaîne de travail fonctionnait de façon chronophage. A l'origine, les experts métiers travaillaient de manière étanche, et toute modification de l'interface publique, une correction de coquille ou une mise en valeur éditoriale nécessitait l'intervention manuelle du service informatique. Le système était perçu comme une \enquote{boîte noire technique} par les experts métiers. Pour décloisonner ces processus, Codex introduit un \enquote{back-offcie collaboratif structuré autour de profils et de rôles spécifiques}. Un système de gestion dees rôles et des droits d'accès, comme l'éditeur, le catalogueur, pour permettre aux équipes métiers d'administrer la plateforme et de corriger des données de manière simplifiée sans dépendre constamment du service informatique. Comme l'a expliqué Thomas Chaineux dans son étude, \enquote{l'application d'administration native Django (\enquote{admin app}) a été historiquement conçue au sein de rédactions de presse pour permettre à des agents non techniciens de saisir et de modifier des données en temps réel de manière autonome et sécurisée}. Cette gestion des rôles permettrait à la bibliothèque Cujas, par exemple, \enquote{à un agent ou à un chercheur de rédiger et publier un texte sur le fonds qu'il a traité, sans pour autant lui donner les droits d'administrateur global sur tout le système}.    
